\section{Prototype evaluation}
\label{sec:experiments}

\subsection{A multi-file sequence pipeline exercises the engineering layer}
The engineering example comprises \EngineeringFiles{} Python source files across \EngineeringCodebases{} codebase roots. Its manifest defines \EngineeringBlocks{} semantic blocks and \EngineeringFlows{} named flows. Membership includes shared numerical and serialization utilities in the second root. The example uses \DatasetRows{} synthetic sequence records; their \DatasetColumns{} columns are summarized directly from the scanned file in Table~\ref{tab:data}. The deterministic generator and its seed are included in the release.

The pipeline prepares arrays, applies a frozen single-block attention encoder, trains a linear regression head, adapts that head, and evaluates held-out records. The attention module performs query, key, and value projections followed by scaled dot products and softmax. The parent block adds residual paths and a feed-forward transformation. Encoder parameters remain fixed throughout. Only the regression head is optimized, so this example supplies no evidence about full Transformer training or model quality on an external task.

The retained run produced a prepared dataset and two checkpoints. It also wrote predictions and loss records. Main-thread events covered \ObservedFiles{} selected source files, and all \ShapeContracts{} declared array-shape contracts matched observed calls. One scalar squared-error contract was bound to the prediction writer's block and passed separately. The architecture audit accepted the current run and its artifacts.

The recorded head-training loss changed from \TrainingInitial{} to \TrainingFinal{} over \TrainingSteps{} updates. Adaptation loss changed from \AdaptationInitial{} to \AdaptationFinal{} over \AdaptationSteps{} updates. The held-out mean squared error (MSE) was \ValidationMSE{} across \ValidationSamples{} records. These values confirm that optimization and evaluation produced inspectable numerical artifacts; they are not a comparison against another learning method.

\begin{table}[t]
\centering
\small
\input{generated/dataset-table}
\caption{The dataset scan found a complete synthetic table and generated its field summary directly from observed records. Ranges describe this fixture and carry no inferred domain meaning.}
\label{tab:data}
\end{table}

\subsection{An authored corpus exercises the declared failure boundary}
The evaluation comprised \CorpusCases{} synthetic cases specified before the reported run. \ValidCases{} represented valid implementations, including an algebraic rewrite and a pseudocode source. The remaining cases exercised formula drift, missing implementations, unsupported control paths, domain failures, and floating-point divergence. Each case retained its exact paper source, implementation, manifest, and generated receipt. Table~\ref{tab:corpus} was generated from those run records.

Each runnable case requested 64 distinct inputs with seed 1729 and tolerance $10^{-10}$. Proof queries had a 3000 millisecond timeout per query. The released run recorded CPython 3.13.12 and Z3 4.16.0 on an ARM64 macOS host. The total number of executed samples was \ExecutedSamples{}. Cases rejected before translation generated no execution evidence. The corpus is a regression suite authored alongside the tool, with no statistical sampling model or real-world prevalence estimate.

\begin{table}[htbp]
\centering
\small
\begin{tabular}{p{.39\linewidth}lll}
\toprule
Authored case & Target & Observed & Real check \\
\midrule
exact square & Pass & Pass & Proved \\
equivalent expansion & Pass & Pass & Proved \\
explicit epsilon & Pass & Pass & Proved \\
pseudocode square & Pass & Pass & Proved \\
constant drift & Fail & Fail & Counterexample \\
missing cross term & Fail & Fail & Counterexample \\
ignored epsilon & Fail & Fail & Counterexample \\
identity for square & Fail & Fail & Counterexample \\
hidden fallback & Blocked & Blocked & Not run \\
commented plan & Blocked & Blocked & Not run \\
discarded computation & Blocked & Blocked & Not run \\
unknown equation & Blocked & Blocked & Not run \\
missing implementation & Blocked & Blocked & Not run \\
division domain hole & Fail & Fail & Domain error \\
vacuous domain & Fail & Fail & Empty domain \\
float cancellation & Fail & Fail & Proved \\
unmodeled call & Blocked & Blocked & Not run \\
pseudocode branch & Blocked & Blocked & Not run \\
\bottomrule
\end{tabular}

\caption{The prototype matched all authored classifications in the released synthetic corpus. Blocked cases lack accepted executable semantics; they are counted separately from demonstrated inequivalence.}
\label{tab:corpus}
\end{table}

\subsection{The evidence layers produced different decisions}
EqTrace matched \ExpectedMatches{} of \CorpusCases{} authored expectations. It accepted \ValidCases{} valid cases and produced no strict pass among the \InvalidCases{} invalid or unsupported cases. These counts describe the provided fixtures. They do not establish sensitivity or specificity on independent research repositories.

Two internal ablations exposed why one check is insufficient. The midpoint-only ablation accepted \MidpointFalsePasses{} invalid cases that the combined check rejected. The ordered-structure ablation rejected \StructureFalseRejects{} valid algebraic rewrite. Both ablations reused EqTrace's strict frontends; unavailable results were not credited as detected bugs. They are component comparisons within this prototype, rather than evaluations of external competing systems.

The floating-point cancellation case passed the real-equivalence query and failed the execution comparison. Its implementation added and subtracted $10^{16}$ around a scalar input. The resulting real expression equals that input, while finite-precision evaluation lost information for tested values. This case demonstrates why \texttt{PROVED\_REAL} and execution success occupy separate fields in the receipt.

\subsection{Self-application checks the arithmetic used by the checker}
The self-application manifest linked Equations~\ref{eq:self-square} and~\ref{eq:self-normalized}, together with Algorithm~\ref{alg:self-normalized}, to functions in the released package. All \SelfContracts{} contracts passed their real-equivalence and sampled execution checks. The normalized discrepancy function is called by the runtime comparison harness. This closes a small paper-to-code-to-evidence cycle within the repository.

The self-application result has a circular trust limitation. The same implementation parses the paper statement, lowers the code, and judges their relationship. A shared translation bug could survive the check. Independent proof-certificate checking, a second frontend implementation, and external fixtures would provide different evidence, and remain future work.

\subsection{Regression and browser checks exercise failure handling}
The released local test run passed \UnitTests{} tests. Engineering regression cases changed source files, datasets, and outputs after a run; supplied a pre-existing unwritten artifact; imported a module without calling its block function; and declared an incorrect array dimension. Each case required a blocking result. The suite also exercised scalar parsing, domain failures, numerical disagreement, stale receipts, and constrained live comparisons.

Browser checks exercised named flows, cross-codebase block inspection, dataset-to-LaTeX display, observed array shapes, and attached scalar proof evidence. Separate checks exercised equation rendering, merged graph inspection, live mismatch detection, exports, and a narrow viewport. These checks establish selected user-interface behavior on Chromium. They do not establish accessibility conformance or behavior in every browser.

The package also scans its own Python roots into semantic blocks and binds its three local self-contracts to the checker block. This view remains a static architecture observation. The release does not label the entire checker as dynamically covered or mathematically verified.
